For compactness in the following proof, write

\[
T_s=\operatorname{Tow}_2(s)\qquad(s\ge1)
\]

\begin{lemma}[Exact tower inversion and floor arithmetic]\label{lem:step-016-tower-arithmetic}
Under the tower and iterated-logarithm definitions in
Section~\ref{sec:setup}, the map
\(s\mapsto T_s\) is strictly increasing and unbounded, with
\(T_s\ge2^s\) for every integer \(s\ge1\). The function
\(u\mapsto\log_2^*u\) is nondecreasing on \((0,\infty)\). Moreover, for
every integer \(r\ge1\),

\[
\log_2^*T_r=r
\tag{1}
\]

and, for every integer \(r\ge2\),

\[
\left\lfloor\log_2(T_r+1)\right\rfloor=T_{r-1}.
\tag{2}
\]
\end{lemma}

\begin{proof}
Because \(T_1=2\) and \(T_{s+1}=2^{T_s}\), induction gives
\(T_s\ge2^s\). Indeed, the claim is true at \(s=1\), and the elementary
induction \(2^s\ge s+1\) gives

\[
T_{s+1}=2^{T_s}\ge2^{2^s}\ge2^{s+1}.
\]

Also \(2^m>m\) for every integer \(m\ge2\), so
\(T_{s+1}>T_s\). Thus the tower is strictly increasing and, from
\(T_s\ge2^s\), unbounded.

For monotonicity of log-star, first note that for \(u>1\),

\[
\log_2^*u=1+\log_2^*(\log_2u).
\tag{3}
\]

If \(0<u\le v\), induction on \(\log_2^*v\), using the monotonicity of
\(\log_2\) and (3), gives \(\log_2^*u\le\log_2^*v\). The base case
\(\log_2^*v=0\) means \(v\le1\), so also \(u\le1\), and both log-star
values are zero. In the induction step, \(u\le1\) is immediate; if
\(u>1\), apply the induction hypothesis to
\(0<\log_2u\le\log_2v\) in (3).

For \(0\le j\le r-1\), repeated use of the tower recursion yields

\[
\log_2^{(j)}T_r=T_{r-j}.
\tag{4}
\]

In particular, the \((r-1)\)-fold iterate equals \(T_1=2>1\), whereas
the \(r\)-fold iterate equals \(\log_2 2=1\). By the minimum in the
definition of log-star, (1) follows exactly.

Finally, put \(q=T_{r-1}\). Then \(T_r=2^q>1\), and hence

\[
2^q<T_r+1<2T_r=2^{q+1}.
\]

Taking base-two logarithms gives

\[
q<\log_2(T_r+1)<q+1.
\]

Since \(q\) is an integer, its floor is exactly \(q=T_{r-1}\), proving
(2). \end{proof}

\begin{proposition}[Exact diagonal structural identities]\label{prop:step-016-structural-diagonal}
Under the primitive diagonal specialization \(r\in\mathbb Z_{\ge2}\),
\(k=r\), \(N=T_r\), Lemma~\ref{lem:step-001-cardinality}, Lemma~\ref{lem:step-001-vc}, Proposition~\ref{prop:step-001-product-ld}, and
Lemma~\ref{lem:step-016-tower-arithmetic},

\[
\operatorname{VC}(C_{r,T_r})=r,
\qquad
\operatorname{LD}(C_{r,T_r})=rT_{r-1},
\qquad
|C_{r,T_r}|=(T_r+1)^r.
\tag{5}
\]

With the setting convention that unadorned \(\log\) is natural,

\[
rT_{r-1}\log2
  <\log|C_{r,T_r}|
  <\frac32rT_{r-1}\log2.
\tag{6}
\]

In particular,
\(\log|C_{r,T_r}|=\Theta(rT_{r-1})\) with universal hidden constants.
\end{proposition}

\begin{proof}
The structural identities give, after the exact substitution
\(k=r,N=T_r\),

\[
\operatorname{VC}(C_{r,T_r})=r,
\]

\[
\operatorname{LD}(C_{r,T_r})
=r\left\lfloor\log_2(T_r+1)\right\rfloor,
\]

and \(|C_{r,T_r}|=(T_r+1)^r\). Equation (2) in
Lemma~\ref{lem:step-016-tower-arithmetic} turns the middle display into
\(rT_{r-1}\), proving every exact identity in (5).

For (6), write \(T_r=2^{T_{r-1}}\). The same strict comparison used for
the floor gives

\[
T_{r-1}\log2
  <\log(T_r+1)
  <(T_{r-1}+1)\log2.
\]

Since \(r\ge2\) implies \(T_{r-1}\ge T_1=2\),
\(T_{r-1}+1\le(3/2)T_{r-1}\). Multiplication by \(r\) and the exact
identity \(\log|C_{r,T_r}|=r\log(T_r+1)\) prove (6). The constants
\(\log2\) and \((3/2)\log2\) are absolute and independent of \(r\).
\end{proof}

\begin{lemma}[Iterated logarithm of the diagonal Littlestone dimension]\label{lem:step-016-ld-logstar}
Under Lemma~\ref{lem:step-016-tower-arithmetic} and
Proposition~\ref{prop:step-016-structural-diagonal}, for every integer
\(r\ge2\),

\[
r-1
\le
\log_2^*\operatorname{LD}(C_{r,T_r})
\le r.
\tag{7}
\]

Consequently,

\[
\frac r2
\le
\log_2^*\operatorname{LD}(C_{r,T_r})
\le r,
\tag{8}
\]

so \(\log_2^*\operatorname{LD}(C_{r,T_r})=\Theta(r)\) with constants
\(1/2\) and \(1\), independent of \(r\).
\end{lemma}

\begin{proof}
Let \(q=T_{r-1}\). Proposition~\ref{prop:step-016-structural-diagonal}
gives

\[
\operatorname{LD}(C_{r,T_r})=rq\ge q=T_{r-1}.
\tag{9}
\]

We next prove the upper comparison \(rq\le T_r=2^q\). At \(r=2\),
\(q=T_1=2\) and both sides equal \(4\). For \(r\ge3\), induction gives
\(T_{r-1}\ge r\): the base is \(T_2=4\ge3\), and
\(T_s\ge s+1\) implies \(T_{s+1}=2^{T_s}\ge2^{s+1}\ge s+2\).
Also, \(2^x\ge x^2\) for every integer \(x\ge4\); this follows by
induction from equality at \(x=4\) and
\(2x^2\ge(x+1)^2\) for \(x\ge3\). Therefore, for \(r\ge3\),

\[
T_r=2^q\ge q^2\ge rq.
\tag{10}
\]

Combining (9)--(10), monotonicity of log-star from
Lemma~\ref{lem:step-016-tower-arithmetic}, and its exact tower inversion
gives

\[
\log_2^*T_{r-1}=r-1
\le\log_2^*(rq)
\le\log_2^*T_r=r,
\]

which is (7). Finally, \(r-1\ge r/2\) for every \(r\ge2\), yielding
(8) and the stated uniform \(\Theta\)-constants. \end{proof}

\begin{proposition}[Tower-diagonal Rate Specialization Bridge]\label{prop:step-016-rate-bridge}
Fix the universal constants
\[
a,c_\delta,\varepsilon_0,\alpha_0,\beta_0,N_0
\]
from Proposition \ref{prop:step-015-exact-closure}, before any candidate is
chosen. Thus \(a=b_*/16\), \(c_\delta=d_*\), \(\varepsilon_0=0.1\), and
\(\alpha_0=\beta_0=2^{-13}\). Define
the finite deterministic index

\[
r_0=\min\{s\in\mathbb Z_{\ge2}:T_s\ge N_0\}.
\tag{11}
\]

Under Assumptions~\ref{assump:candidate-regime},
\ref{assump:central-dp}, and
\ref{assump:distribution-free-realizable-pac}, every diagonal candidate
with integer \(r\ge r_0\), \(k=r\), and \(N=T_r\) satisfies

\[
n\ge ar^2
\tag{12}
\]

and

\[
n\ge
a\,\operatorname{VC}(C_{r,T_r})
   \log_2^*\operatorname{LD}(C_{r,T_r}).
\tag{13}
\]

Moreover,

\[
\frac12r^2
\le
\operatorname{VC}(C_{r,T_r})
\log_2^*\operatorname{LD}(C_{r,T_r})
\le r^2.
\tag{14}
\]

Thus (12)--(13) are exactly
\(n=\Omega(r^2)=\Omega(\operatorname{VC}\log_2^*\operatorname{LD})\),
with constants independent of \(r\), and only within the approved
candidate/privacy/PAC regime.
\end{proposition}

\begin{proof}
Lemma~\ref{lem:step-016-tower-arithmetic} proves that \(T_s\) is
increasing and unbounded, so the set in (11) is nonempty and \(r_0\) is a
finite integer depending only on the already-fixed \(N_0\). For every
\(r\ge r_0\), one has \(r\ge2\) and \(T_r\ge T_{r_0}\ge N_0\).
Therefore \(k=r,N=T_r\) meet the structural range of Proposition~\ref{prop:step-015-exact-closure}.

No privacy condition is absorbed in this substitution. The candidate has
the exact budget

\[
m_{n,r}=\max\left\{8,\left\lceil\frac{4n}{r}\right\rceil\right\},
\]

and Assumption~\ref{assump:candidate-regime} remains exactly

\[
n\in\mathbb Z_{\ge1},
\quad 0<\varepsilon\le\varepsilon_0,
\quad
0<\delta\le
\min\left\{
\frac1{n\log(n+1)},
\frac{c_\delta}{m_{n,r}^2\log(m_{n,r}+1)}
\right\}.
\tag{15}
\]

The learner remains an arbitrary randomized map

\[
A:(X_{r,T_r}\times\{0,1\})^n
  \longrightarrow\{0,1\}^{X_{r,T_r}}
\]

satisfying the exact central replacement-DP and distribution-free
realizable-PAC assumptions at this same candidate. Proposition~\ref{prop:step-015-exact-closure} therefore applies without a
mode or scope conversion and gives

\[
n\ge ar\log_2^*T_r.
\]

Equation (1) in Lemma~\ref{lem:step-016-tower-arithmetic} makes the
right-hand side exactly \(ar^2\), proving (12). No asymptotic term is
dropped in this equality.

By Proposition~\ref{prop:step-016-structural-diagonal} and
Lemma~\ref{lem:step-016-ld-logstar},

\[
\operatorname{VC}(C_{r,T_r})=r,
\qquad
\frac r2
\le\log_2^*\operatorname{LD}(C_{r,T_r})\le r.
\]

Multiplication proves (14). In particular, the upper half of (14) and
(12) give

\[
n\ge ar^2
\ge
a\,\operatorname{VC}(C_{r,T_r})
   \log_2^*\operatorname{LD}(C_{r,T_r}),
\]

which is (13) with the same inherited constant \(a\).

The finite initial indices are not hidden inside an \(r\)-dependent
constant. At \(r=2\), the arithmetic itself is exact:

\[
T_1=2,\qquad T_2=4,\qquad \log_2^*T_2=2,\qquad
\lfloor\log_2(5)\rfloor=2,
\]

and hence \(\operatorname{VC}=2\), \(\operatorname{LD}=4\),
\(|C|=25\), and \(\log_2^*\operatorname{LD}=2\). The same structural
identities hold for every \(2\le r<r_0\). By the minimality of \(r_0\)
and monotonicity of the tower, every such index has \(T_r<N_0\). Therefore
Proposition~\ref{prop:step-015-exact-closure} is inapplicable, so no learning
lower-bound conclusion is asserted for those indices. The asymptotic statements begin
at the single fixed index \(r_0\). \end{proof}

\begin{proposition}[Diagonal separation and unresolved scales]\label{prop:step-016-scale-comparison}
Under Proposition~\ref{prop:step-016-structural-diagonal} and
Lemma~\ref{lem:step-016-ld-logstar}, for every integer \(r\ge2\),

\[
2r-1
\le
\operatorname{VC}(C_{r,T_r})
+\log_2^*\operatorname{LD}(C_{r,T_r})
\le2r,
\tag{16}
\]

and

\[
r2^{r-1}
\le
\operatorname{VC}(C_{r,T_r})
2^{\log_2^*\operatorname{LD}(C_{r,T_r})}
\le r2^r.
\tag{17}
\]

Consequently, the established lower-bound scale \(r^2\) exceeds the
additive expression by an unbounded factor, but

\[
r^2
=o\!\left(
\operatorname{VC}(C_{r,T_r})
2^{\log_2^*\operatorname{LD}(C_{r,T_r})}
\right)
\tag{18}
\]

and

\[
r^2=o(\log|C_{r,T_r}|).
\tag{19}
\]

These comparisons concern the size of the proved lower-bound expression;
they are not upper bounds on the true sample complexity.
\end{proposition}

\begin{proof}
Substitute \(\operatorname{VC}=r\) and the two bounds in (7). Addition
gives (16), and exponentiation followed by multiplication by \(r\) gives
(17). In particular,

\[
\frac{r^2}{
 \operatorname{VC}(C_{r,T_r})
 +\log_2^*\operatorname{LD}(C_{r,T_r})}
\ge\frac r2\longrightarrow\infty,
\]

which proves the unbounded-factor comparison with the additive scale.

For the first remaining scale, the lower half of (17) gives

\[
0\le
\frac{r^2}{
 \operatorname{VC}(C_{r,T_r})
 2^{\log_2^*\operatorname{LD}(C_{r,T_r})}}
\le\frac r{2^{r-1}}.
\tag{20}
\]

The right-hand side tends to zero: its ratio at consecutive indices is
\((r+1)/(2r)\le3/4\) for \(r\ge2\). This proves (18).

For cardinality, (6) and \(T_{r-1}\ge2^{r-1}\) from
Lemma~\ref{lem:step-016-tower-arithmetic} yield

\[
0\le\frac{r^2}{\log|C_{r,T_r}|}
<\frac{r}{T_{r-1}\log2}
\le\frac{r}{2^{r-1}\log2}\longrightarrow0,
\]

which proves (19). \end{proof}
